% =============================================================================
%  entretien_python_backend.tex
%  Preparation entretien technique -- Backend Python (Django / FastAPI)
%  300+ questions/reponses classees par niveau (Junior -> Mid -> Senior)
%
%  Compilation (deux passes pour la table des matieres) :
%     xelatex entretien_python_backend.tex
%     xelatex entretien_python_backend.tex
%  (ou pdflatex entretien_python_backend.tex, deux fois)
% =============================================================================
\documentclass[11pt,a4paper]{article}

\usepackage[utf8]{inputenc}
\usepackage[T1]{fontenc}
\usepackage[french]{babel}
\usepackage[margin=2.2cm]{geometry}
\usepackage[table]{xcolor}
\usepackage{listings}
\usepackage{tcolorbox}
\usepackage{enumitem}
\usepackage{titlesec}
\usepackage{longtable}
\usepackage{array}
\usepackage{fancyhdr}
\usepackage{hyperref}

% --- Couleurs ---------------------------------------------------------------
\definecolor{primary}{HTML}{1F3A5F}
\definecolor{accent}{HTML}{2E86DE}
\definecolor{codebg}{HTML}{F4F7FB}
\definecolor{codekw}{HTML}{1F3A5F}
\definecolor{codestr}{HTML}{27AE60}
\definecolor{codecom}{HTML}{7F8C99}
\definecolor{rulecol}{HTML}{D5DEE8}
\definecolor{junior}{HTML}{27AE60}
\definecolor{mid}{HTML}{E67E22}
\definecolor{senior}{HTML}{C0392B}
\definecolor{qbg}{HTML}{EAF1FB}

% --- Liens ------------------------------------------------------------------
\hypersetup{colorlinks=true, linkcolor=accent, urlcolor=accent,
  pdftitle={Entretien Technique - Backend Python (Django/FastAPI)}}

% --- En-tete / pied de page -------------------------------------------------
\pagestyle{fancy}
\fancyhf{}
\fancyhead[L]{\small\color{primary}Entretien Backend Python -- Django / FastAPI}
\fancyhead[R]{\small\color{codecom}\thepage}
\renewcommand{\headrulewidth}{0.4pt}

% --- Titres de section ------------------------------------------------------
\titleformat{\section}{\Large\bfseries\color{primary}}{\thesection.}{0.6em}{}
\titleformat{\subsection}{\large\bfseries\color{accent}}{}{0em}{}
\titleformat{\subsubsection}{\normalsize\bfseries\color{primary}}{}{0em}{}

% --- Style du code Python ---------------------------------------------------
\lstdefinestyle{py}{
  language=Python,
  backgroundcolor=\color{codebg},
  basicstyle=\ttfamily\small,
  keywordstyle=\color{codekw}\bfseries,
  stringstyle=\color{codestr},
  commentstyle=\color{codecom}\itshape,
  numbers=none,
  frame=single,
  rulecolor=\color{rulecol},
  breaklines=true,
  showstringspaces=false,
  columns=fullflexible,
  keepspaces=true,
  aboveskip=6pt, belowskip=6pt,
  xleftmargin=6pt, xrightmargin=6pt,
  literate=%
    {é}{{\'e}}1 {è}{{\`e}}1 {ê}{{\^e}}1 {à}{{\`a}}1 {ô}{{\^o}}1
    {î}{{\^i}}1 {ç}{{\c c}}1 {É}{{\'E}}1 {È}{{\`E}}1 {ù}{{\`u}}1,
}
\lstset{style=py}

% --- Boites "Astuce" / "Piege" -----------------------------------------------
\newtcolorbox{infobox}[1]{colback=codebg,colframe=accent,title=#1,
  fonttitle=\bfseries\color{white},boxrule=0.5pt,arc=2pt,left=6pt,right=6pt,top=4pt,bottom=4pt}

\newtcolorbox{piege}[1][Piege frequent]{colback=red!4,colframe=senior,title=#1,
  fonttitle=\bfseries\color{white},boxrule=0.5pt,arc=2pt,left=6pt,right=6pt,top=4pt,bottom=4pt}

% --- Badges de niveau ---------------------------------------------------------
\newcommand{\badgeJ}{\textcolor{white}{\colorbox{junior}{\small\bfseries\ JUNIOR\ }}}
\newcommand{\badgeM}{\textcolor{white}{\colorbox{mid}{\small\bfseries\ MID\ }}}
\newcommand{\badgeS}{\textcolor{white}{\colorbox{senior}{\small\bfseries\ SENIOR\ }}}

% --- Compteur de questions global --------------------------------------------
\newcounter{qcount}

% --- Environnement Question/Reponse ------------------------------------------
% Usage: \q{NIVEAU}{Texte de la question}
%        Reponse ...
%        \finq
\newenvironment{qa}[2]{%
  \stepcounter{qcount}%
  \par\vspace{6pt}%
  \noindent\colorbox{qbg}{\parbox{\dimexpr\linewidth-12pt}{\vspace{2pt}%
  \ifnum#1=1 \badgeJ\fi\ifnum#1=2 \badgeM\fi\ifnum#1=3 \badgeS\fi
  \ \textbf{Q\theqcount.} #2\vspace{2pt}}}\par\vspace{4pt}%
  \noindent\textbf{\color{accent}Reponse.}\ \ignorespaces
}{%
  \par\vspace{8pt}%
}

% =============================================================================
\begin{document}

% --- Page de titre ----------------------------------------------------------
\begin{titlepage}
\centering
\vspace*{2.5cm}
{\color{primary}\Huge\bfseries Entretien Technique\par}
\vspace{0.4cm}
{\color{primary}\Huge\bfseries Backend Python\par}
\vspace{0.6cm}
{\color{accent}\LARGE Django \ $\cdot$\ FastAPI \ $\cdot$\ DRF \ $\cdot$\ PostgreSQL \ $\cdot$\ Tests \ $\cdot$\ Architecture\par}
\vspace{1.2cm}
{\large Plus de 300 questions--reponses classees\\
Junior \ $\rightarrow$\ Mid \ $\rightarrow$\ Senior\par}
\vspace{1cm}

\textcolor{white}{\colorbox{junior}{\small\bfseries\ JUNIOR\ }}\quad
\textcolor{white}{\colorbox{mid}{\small\bfseries\ MID\ }}\quad
\textcolor{white}{\colorbox{senior}{\small\bfseries\ SENIOR\ }}

\vfill
{\color{codecom}\large Support de revision a imprimer\par}
\vspace{2cm}
\end{titlepage}

\tableofcontents
\newpage

% =============================================================================
\section{Comment utiliser ce document}
% =============================================================================
Ce manuel couvre les questions qui reviennent le plus souvent en entretien pour
un poste \textbf{Backend Software Engineer Python} (Django et/ou FastAPI), sur
la base d'une synthese de sources 2025--2026 (TechPrep, Learnixo, HackerX,
Medium, Reddit, Baeldung).

\begin{infobox}{Repartition observee d'un entretien type}
\textbf{Python core} $\approx$ 35\,\% \quad $\vert$ \quad
\textbf{Framework (Django/FastAPI)} $\approx$ 25\,\% \quad $\vert$ \quad
\textbf{Base de donnees / SQL} $\approx$ 15\,\% \quad $\vert$ \quad
\textbf{Tests, securite, archi, deploiement} $\approx$ 10\,\% \quad $\vert$ \quad
\textbf{Questions sur tes projets} $\approx$ 30$-$40\,\%
\end{infobox}

\textbf{Niveaux.} Chaque question porte un badge :
\badgeJ\ (bases attendues de tous), \badgeM\ (comprehension pratique /
production), \badgeS\ (architecture, performance, decisions de conception).

\textbf{Conseil de revision.} Pour un profil deja experimente (FastAPI, RAG,
Kafka, Spark, PostgreSQL, microservices), concentre-toi en priorite sur les
sections \emph{Async Python}, \emph{FastAPI en production}, \emph{Django ORM \&
performance}, \emph{PostgreSQL}, \emph{Securite} et \emph{Questions projets} --
les bases Python doivent rester des reflexes, pas une revision.

\newpage

% =============================================================================
\section{Python Core}
% =============================================================================

\subsection{Types de donnees}

\begin{qa}{1}{Quelle est la difference entre une \texttt{list} et un \texttt{tuple} ?}
Les deux sont des sequences ordonnees, mais la \texttt{list} est \emph{mutable}
(on peut ajouter/retirer/modifier des elements) alors que le \texttt{tuple} est
\emph{immutable} (sa taille et son contenu sont fixes apres creation). Consequence
pratique : un tuple est hashable (si ses elements le sont) et peut donc servir
de cle de dictionnaire ou d'element d'un \texttt{set}, contrairement a une liste.
Les tuples sont aussi legerement plus rapides a creer et consomment moins de
memoire.
\end{qa}

\begin{qa}{1}{Pourquoi un tuple est-il immutable ? Quel interet concret ?}
L'immutabilite est une decision de conception : elle garantit qu'une structure
ne changera pas apres sa creation, ce qui la rend sure a partager entre
plusieurs parties d'un programme (pas d'effet de bord cache) et hashable (donc
utilisable comme cle de \texttt{dict} ou element de \texttt{set}). C'est aussi
la base du \emph{unpacking} multiple valeurs de retour (\texttt{return x, y}).
\end{qa}

\begin{qa}{1}{Quand utiliser un \texttt{set} plutot qu'une \texttt{list} ?}
Des que l'on a besoin (1) de tester l'appartenance frequemment -- \texttt{in}
est en \texttt{O(1)} moyen sur un \texttt{set} contre \texttt{O(n)} sur une
liste -- ou (2) d'eliminer des doublons, ou (3) de faire des operations
ensemblistes (union \texttt{|}, intersection \texttt{\&}, difference
\texttt{-}). Contrepartie : un \texttt{set} n'est pas ordonne et ses elements
doivent etre hashables.
\end{qa}

\begin{qa}{2}{Que signifie \og{}hashable\fg{} et pourquoi les cles de dict doivent-elles l'etre ?}
Un objet est hashable s'il possede une methode \texttt{\_\_hash\_\_} qui renvoie
une valeur entiere stable pendant toute sa duree de vie, utilisee par le
dictionnaire pour placer/retrouver la cle en \texttt{O(1)} moyen dans sa table
de hachage. Si l'objet est mutable (comme une liste), son hash pourrait changer
apres modification, ce qui casserait la structure interne du dict -- c'est
pourquoi les listes ne sont pas hashables mais les tuples le sont (si leur
contenu l'est aussi).
\end{qa}

\begin{qa}{2}{Difference entre \texttt{set} et \texttt{frozenset} ?}
\texttt{frozenset} est la version immutable de \texttt{set} : une fois cree, on
ne peut plus y ajouter/retirer d'elements. Il est donc hashable et peut servir
de cle de dictionnaire ou d'element d'un autre set, ce qu'un \texttt{set}
classique ne permet pas.
\end{qa}

\subsection{Variables, mutabilite, copies}

\begin{qa}{1}{Pourquoi une liste passee en argument de fonction peut-elle etre modifiee par la fonction ?}
En Python, les arguments sont passes \og{}par reference d'objet\fg{}
(\emph{pass by object reference}) : la variable locale de la fonction pointe
vers le \emph{meme} objet que l'appelant. Si cet objet est mutable (liste,
dict, set) et que la fonction le modifie \emph{en place} (\texttt{.append()},
\texttt{[i]=...}), le changement est visible a l'exterieur. En revanche,
reassigner le parametre (\texttt{param = []}) ne fait que rebrancher la
variable locale sur un nouvel objet, sans effet sur l'appelant.
\begin{lstlisting}
def ajoute(l):
    l.append(4)      # modifie l'objet partage -> visible dehors

def remplace(l):
    l = [9, 9, 9]     # rebranche la variable locale -> invisible dehors
\end{lstlisting}
\end{qa}

\begin{qa}{2}{Difference entre \emph{shallow copy} et \emph{deep copy} ?}
Une copie superficielle (\texttt{copy.copy()}, ou \texttt{list(l)}, \texttt{l[:]})
cree un nouvel objet conteneur mais dont les elements internes restent des
\emph{references} vers les memes objets que l'original -- si un element est lui
meme mutable (liste imbriquee), le modifier affecte les deux copies. Une copie
profonde (\texttt{copy.deepcopy()}) reconstruit recursivement tous les objets
imbriques, donc les deux structures deviennent totalement independantes.
\begin{lstlisting}
import copy
a = [[1, 2], [3, 4]]
b = copy.copy(a)       # shallow
c = copy.deepcopy(a)   # deep
b[0].append(99)
print(a)  # [[1, 2, 99], [3, 4]]  -> a est impacte (meme sous-liste)
print(c)  # [[1, 2], [3, 4]]      -> c est independant
\end{lstlisting}
\end{qa}

\begin{qa}{1}{Quelle est la difference entre \texttt{==} et \texttt{is} ?}
\texttt{==} compare l'\emph{egalite de valeur} (appelle \texttt{\_\_eq\_\_}) ;
\texttt{is} compare l'\emph{identite} -- est-ce le meme objet en memoire
(memes \texttt{id()}) ? Deux listes contenant les memes elements sont
\texttt{==} mais pas \texttt{is}. On utilise \texttt{is} typiquement pour
comparer a \texttt{None}, \texttt{True}/\texttt{False} (\texttt{if x is None}),
car ce sont des singletons.
\end{qa}

\begin{qa}{2}{Pourquoi \texttt{a is b} peut etre \texttt{True} pour deux petits entiers identiques mais \texttt{False} pour de grands entiers ?}
CPython met en cache (\emph{interning}) les petits entiers entre $-5$ et $256$
ainsi que certaines chaines courtes : deux references vers la valeur \texttt{5}
pointent donc vers le meme objet. Au-dela de cette plage, chaque litteral entier
cree potentiellement un nouvel objet, donc \texttt{is} peut renvoyer
\texttt{False} meme si les valeurs sont egales. C'est un detail d'implementation
qu'il ne faut jamais utiliser comme mecanisme de comparaison volontaire --
toujours utiliser \texttt{==} pour comparer des valeurs.
\end{qa}

\subsection{Fonctions}

\begin{qa}{1}{A quoi servent \texttt{*args} et \texttt{**kwargs} ?}
\texttt{*args} collecte un nombre variable d'arguments positionnels dans un
\texttt{tuple} ; \texttt{**kwargs} collecte les arguments nommes restants dans
un \texttt{dict}. Ils permettent d'ecrire des fonctions a signature flexible
(wrappers, decorateurs, fonctions generiques) sans connaitre a l'avance le
nombre de parametres.
\begin{lstlisting}
def f(*args, **kwargs):
    print(args)    # (1, 2, 3)
    print(kwargs)  # {'x': 10}

f(1, 2, 3, x=10)
\end{lstlisting}
\end{qa}

\begin{qa}{1}{Difference entre argument positionnel et \emph{keyword argument} ?}
Un argument positionnel est associe au parametre selon sa position dans
l'appel ; un \emph{keyword argument} est passe explicitement par
\texttt{nom=valeur}, ce qui rend l'appel plus lisible et permet de sauter des
parametres ayant une valeur par defaut, quel que soit leur ordre de
declaration (tant qu'ils viennent apres les arguments positionnels).
\end{qa}

\begin{qa}{2}{Quel est le piege classique des arguments par defaut mutables ?}
Les valeurs par defaut sont evaluees \textbf{une seule fois}, a la definition
de la fonction, et partagees entre tous les appels. Utiliser une liste ou un
dict mutable comme valeur par defaut cree donc un etat partage et persistant
entre appels -- un bug tres classique.
\begin{lstlisting}
def ajoute(item, liste=[]):   # BUG : meme liste reutilisee a chaque appel
    liste.append(item)
    return liste

def ajoute_ok(item, liste=None):
    if liste is None:
        liste = []
    liste.append(item)
    return liste
\end{lstlisting}
\begin{piege}
\texttt{ajoute(1)} puis \texttt{ajoute(2)} renvoie \texttt{[1, 2]} et non
\texttt{[2]} : c'est la meme liste par defaut qui s'accumule silencieusement.
\end{piege}
\end{qa}

\subsection{Programmation orientee objet}

\begin{qa}{1}{Rappelle les quatre piliers de l'OOP.}
\textbf{Encapsulation} : regrouper donnees et methodes, limiter l'acces direct
aux details internes (convention \texttt{\_protected} / \texttt{\_\_private}).
\textbf{Heritage} : une classe fille reutilise/etend le comportement d'une
classe mere. \textbf{Polymorphisme} : des objets de classes differentes
repondent au meme appel de methode chacun a sa maniere (duck typing en
Python). \textbf{Abstraction} : exposer une interface simple en masquant la
complexite d'implementation (classes abstraites, \texttt{ABC}).
\end{qa}

\begin{qa}{1}{Quelle difference entre \texttt{@staticmethod} et \texttt{@classmethod} ?}
Une \texttt{@classmethod} recoit automatiquement la classe en premier
argument (\texttt{cls}) et peut donc acceder/modifier l'etat de classe ou
servir de constructeur alternatif. Une \texttt{@staticmethod} ne recoit ni
\texttt{self} ni \texttt{cls} : c'est une fonction normale simplement
rattachee a la classe par organisation logique, sans acces automatique a
l'instance ou a la classe.
\begin{lstlisting}
class Pizza:
    def __init__(self, ingredients):
        self.ingredients = ingredients

    @classmethod
    def margherita(cls):          # constructeur alternatif
        return cls(["tomate", "mozzarella"])

    @staticmethod
    def temps_cuisson(taille_cm):  # utilitaire sans lien avec l'instance
        return taille_cm * 0.5
\end{lstlisting}
\end{qa}

\begin{qa}{1}{Difference entre \emph{instance method}, \texttt{classmethod} et \texttt{staticmethod} ?}
La methode d'instance (par defaut, premier parametre \texttt{self}) agit sur
une instance precise et peut lire/modifier son etat. La \texttt{classmethod}
agit sur la classe elle-meme (\texttt{cls}), utile pour les constructeurs
alternatifs ou l'etat partage entre instances. La \texttt{staticmethod} n'agit
ni sur l'un ni sur l'autre : c'est un simple regroupement logique.
\end{qa}

\begin{qa}{2}{Comment fonctionne l'heritage multiple et le MRO (Method Resolution Order) ?}
Python utilise l'algorithme C3 linearization pour determiner l'ordre dans
lequel les classes parentes sont consultees lors de la resolution d'une
methode (\texttt{ClassName.\_\_mro\_\_} ou \texttt{help(ClassName)}). Cela
garantit un ordre coherent meme avec des heritages en diamant, et c'est ce qui
permet a \texttt{super()} de fonctionner correctement en cooperant avec
plusieurs classes parentes.
\end{qa}

\begin{qa}{2}{Qu'est-ce que le \emph{duck typing} ?}
Principe selon lequel le type d'un objet importe moins que les methodes qu'il
implemente : \og{}si ca marche comme un canard et cancane comme un canard,
c'est un canard\fg{}. Python ne verifie pas de type explicite a la compilation
: si un objet possede la methode attendue (par ex. \texttt{.read()}), il peut
etre utilise la ou un fichier est attendu, meme sans heriter d'une classe
commune.
\end{qa}

\subsection{Methodes speciales (dunder)}

\begin{qa}{1}{Difference entre \texttt{\_\_init\_\_} et \texttt{\_\_new\_\_} ?}
\texttt{\_\_new\_\_} est la methode statique qui \emph{cree} et renvoie une
nouvelle instance (allocation memoire) ; \texttt{\_\_init\_\_} est appelee
\emph{apres}, sur l'instance deja creee, pour l'\emph{initialiser} (definir ses
attributs). On surcharge \texttt{\_\_new\_\_} dans des cas rares : classes
immutables (heritage de \texttt{tuple}, \texttt{str}), singletons, ou
metaclasses.
\end{qa}

\begin{qa}{1}{Difference entre \texttt{\_\_str\_\_} et \texttt{\_\_repr\_\_} ?}
\texttt{\_\_str\_\_} definit la representation \og{}lisible\fg{} destinee a
l'utilisateur final (appelee par \texttt{print()} et \texttt{str()}).
\texttt{\_\_repr\_\_} definit une representation \og{}non ambigue\fg{}
destinee au developpeur/debug (appelee dans le REPL, par \texttt{repr()}, et en
secours si \texttt{\_\_str\_\_} n'est pas defini). Regle courante :
\texttt{\_\_repr\_\_} devrait, si possible, ressembler a du code Python valide
permettant de recreer l'objet.
\begin{lstlisting}
class Point:
    def __init__(self, x, y):
        self.x, self.y = x, y
    def __repr__(self):
        return f"Point(x={self.x}, y={self.y})"
    def __str__(self):
        return f"({self.x}, {self.y})"
\end{lstlisting}
\end{qa}

\begin{qa}{2}{Cite d'autres dunder methods utiles et leur role.}
\texttt{\_\_len\_\_} (pour \texttt{len(obj)}), \texttt{\_\_eq\_\_}/\texttt{\_\_lt\_\_}
(comparaisons), \texttt{\_\_iter\_\_}/\texttt{\_\_next\_\_} (iteration),
\texttt{\_\_getitem\_\_}/\texttt{\_\_setitem\_\_} (indexation \texttt{obj[i]}),
\texttt{\_\_enter\_\_}/\texttt{\_\_exit\_\_} (context manager \texttt{with}),
\texttt{\_\_call\_\_} (rendre l'objet appelable comme une fonction),
\texttt{\_\_hash\_\_} (utilisation comme cle de dict).
\end{qa}

\subsection{Decorateurs}

\begin{qa}{1}{Qu'est-ce qu'un decorateur ?}
Une fonction qui prend une autre fonction (ou classe) en argument et renvoie
une version modifiee/enrichie de celle-ci, sans changer son code source. Cela
repose sur le fait que les fonctions sont des objets de premiere classe en
Python (on peut les passer en argument, les stocker, les renvoyer).
\begin{lstlisting}
def chrono(fonction):
    def wrapper(*args, **kwargs):
        debut = time.time()
        resultat = fonction(*args, **kwargs)
        print(f"{fonction.__name__} : {time.time() - debut:.3f}s")
        return resultat
    return wrapper

@chrono
def calcul_lourd():
    ...
\end{lstlisting}
\end{qa}

\begin{qa}{2}{Comment fonctionne \texttt{@property} ?}
\texttt{@property} transforme une methode en attribut en lecture (accessible
sans parentheses), permettant d'ajouter de la logique (validation, calcul) tout
en gardant une API simple type attribut. Combinee a \texttt{@x.setter}, elle
permet de controler l'ecriture (par ex. valider une valeur avant de
l'assigner) sans casser l'interface publique de la classe.
\begin{lstlisting}
class Compte:
    def __init__(self, solde):
        self._solde = solde

    @property
    def solde(self):
        return self._solde

    @solde.setter
    def solde(self, valeur):
        if valeur < 0:
            raise ValueError("Solde negatif interdit")
        self._solde = valeur
\end{lstlisting}
\end{qa}

\begin{qa}{2}{Comment ecrire un decorateur qui accepte des parametres ?}
Il faut un niveau d'imbrication supplementaire : une fonction qui prend les
parametres du decorateur et renvoie le \og{}vrai\fg{} decorateur.
\begin{lstlisting}
def repeter(n):
    def decorateur(fonction):
        def wrapper(*args, **kwargs):
            for _ in range(n):
                resultat = fonction(*args, **kwargs)
            return resultat
        return wrapper
    return decorateur

@repeter(3)
def dire_bonjour():
    print("Bonjour")
\end{lstlisting}
\end{qa}

\begin{qa}{3}{Pourquoi utiliser \texttt{functools.wraps} dans un decorateur ?}
Sans \texttt{@functools.wraps(fonction)} sur le \texttt{wrapper}, la fonction
decoree perd ses metadonnees d'origine (\texttt{\_\_name\_\_},
\texttt{\_\_doc\_\_}, signature), ce qui casse l'introspection, la
documentation automatique et certains outils (debogueurs, frameworks comme
FastAPI qui inspectent la signature pour l'injection de dependances).
\end{qa}

\subsection{Generateurs et iterateurs}

\begin{qa}{1}{Difference entre un \emph{iterable} et un \emph{iterateur} ?}
Un \emph{iterable} est un objet capable de produire un iterateur (il implemente
\texttt{\_\_iter\_\_}), comme une liste. Un \emph{iterateur} est l'objet qui
garde l'etat de progression et produit les valeurs une a une via
\texttt{\_\_next\_\_}, jusqu'a lever \texttt{StopIteration}. Une liste est
iterable mais n'est pas elle-meme un iterateur : \texttt{iter(liste)} en
produit un nouveau a chaque appel.
\end{qa}

\begin{qa}{2}{A quoi sert \texttt{yield} et comment fonctionne un generateur ?}
\texttt{yield} transforme une fonction en \emph{generateur} : au lieu de
calculer et retourner tout le resultat d'un coup, la fonction produit les
valeurs une par une et \emph{suspend son execution} entre chaque appel,
conservant son etat local. Avantage majeur : la memoire consommee reste
constante meme sur une sequence tres grande (voire infinie), car rien n'est
stocke en liste intermediaire.
\begin{lstlisting}
def carres(n):
    for i in range(n):
        yield i * i

for c in carres(1_000_000):   # pas de liste de 1M elements en memoire
    ...
\end{lstlisting}
\end{qa}

\begin{qa}{2}{Quand preferer un generateur a une liste ?}
Des que la sequence est grande ou potentiellement infinie, ou que l'appelant
n'a pas besoin de tous les elements en meme temps (traitement en flux,
pipeline). La contrepartie : un generateur ne se parcourt qu'une seule fois et
n'offre pas d'acces aleatoire (\texttt{gen[3]} impossible) ni \texttt{len()}.
\end{qa}

\subsection{Gestion memoire}

\begin{qa}{2}{Comment fonctionne le \emph{reference counting} en CPython ?}
Chaque objet garde un compteur du nombre de references qui pointent vers lui.
Quand ce compteur tombe a zero (plus personne ne le reference), l'objet est
immediatement desalloue. C'est le mecanisme principal de gestion memoire de
CPython, complete par un \emph{garbage collector} pour les cas que le comptage
seul ne resout pas.
\end{qa}

\begin{qa}{2}{Pourquoi a-t-on besoin d'un Garbage Collector en plus du reference counting ?}
Le comptage de references seul ne detecte pas les \emph{cycles de reference}
(deux objets qui se referencent mutuellement, meme si plus rien d'exterieur ne
les reference -- leur compteur ne retombe jamais a zero). Le module
\texttt{gc} implemente un algorithme de detection de cycles (generational
garbage collection) qui parcourt periodiquement les objets pour liberer ces
cycles orphelins.
\end{qa}

\begin{qa}{2}{Qu'est-ce que la regle LEGB ?}
Ordre de resolution des noms de variables en Python : \textbf{L}ocal (dans la
fonction courante) $\rightarrow$ \textbf{E}nclosing (fonction englobante, pour
les closures) $\rightarrow$ \textbf{G}lobal (module) $\rightarrow$
\textbf{B}uilt-in (fonctions natives comme \texttt{len}, \texttt{print}).
Python cherche le nom dans cet ordre et s'arrete a la premiere correspondance.
\end{qa}

\begin{qa}{3}{Qu'est-ce que le GIL (Global Interpreter Lock) et quel est son impact ?}
Un verrou global qui n'autorise qu'un seul thread a executer du bytecode
Python a la fois dans un processus CPython, meme sur une machine
multi-coeurs. Consequence : le multithreading n'accelere pas les taches
\emph{CPU-bound} (calcul pur), mais reste efficace pour les taches
\emph{I/O-bound} (le GIL est relache pendant les attentes reseau/disque). Pour
paralleliser du calcul, on utilise \texttt{multiprocessing} (processus
separes, chacun avec son propre GIL) plutot que \texttt{threading}.
\end{qa}

\subsection{Exceptions}

\begin{qa}{1}{Explique le fonctionnement de \texttt{try / except / else / finally}.}
Le bloc \texttt{try} contient le code a risque. \texttt{except} capture et
traite une exception specifique. \texttt{else} s'execute uniquement si
\emph{aucune} exception n'a ete levee. \texttt{finally} s'execute
\emph{toujours}, qu'il y ait eu exception ou non (utile pour liberer des
ressources : fermer un fichier, une connexion).
\begin{lstlisting}
try:
    valeur = int(entree)
except ValueError:
    print("Entree invalide")
else:
    print(f"Conversion reussie : {valeur}")
finally:
    print("Nettoyage effectue")
\end{lstlisting}
\end{qa}

\begin{qa}{2}{Comment creer et utiliser une exception personnalisee ?}
On herite de \texttt{Exception} (ou d'une sous-classe plus specifique) pour
creer un type d'erreur metier explicite, plus facile a capturer et documenter
qu'une exception generique.
\begin{lstlisting}
class SoldeInsuffisantError(Exception):
    def __init__(self, solde, montant):
        self.solde = solde
        self.montant = montant
        super().__init__(
            f"Solde {solde} insuffisant pour retirer {montant}")

def retirer(solde, montant):
    if montant > solde:
        raise SoldeInsuffisantError(solde, montant)
\end{lstlisting}
\end{qa}

\begin{qa}{2}{Quelle difference entre \texttt{raise Erreur()} et \texttt{raise Erreur() from autre\_erreur} ?}
\texttt{raise ... from ...} chaine explicitement une nouvelle exception a sa
cause d'origine, en preservant la traceback complete dans l'attribut
\texttt{\_\_cause\_\_}. Tres utile pour transformer une erreur bas niveau
(ex. \texttt{KeyError}) en une erreur metier explicite sans perdre le contexte
de debogage.
\end{qa}

\subsection{Async Python}

\begin{qa}{2}{Que signifie \texttt{async}/\texttt{await} et qu'est-ce qu'une coroutine ?}
\texttt{async def} definit une \emph{coroutine} : une fonction dont
l'execution peut etre suspendue et reprise plus tard sans bloquer le thread.
\texttt{await} suspend la coroutine courante jusqu'a ce que l'operation
attendue (typiquement une I/O) se termine, en rendant la main a l'\emph{event
loop} pour qu'il execute d'autres taches pendant l'attente.
\end{qa}

\begin{qa}{2}{Qu'est-ce que l'\emph{event loop} dans \texttt{asyncio} ?}
Le coeur du modele asynchrone : une boucle qui gere un ensemble de taches
(coroutines) et bascule de l'une a l'autre chaque fois qu'une tache
\texttt{await} une operation I/O, au lieu de rester bloquee a attendre.
Contrairement au multithreading, tout se passe sur \textbf{un seul thread} :
c'est de la concurrence cooperative, pas du parallelisme.
\end{qa}

\begin{qa}{2}{Pourquoi \texttt{async} peut rendre une application plus rapide alors que Python reste mono-thread (a cause du GIL) ?}
Le gain ne vient pas de calcul parallele mais du fait de ne plus \emph{bloquer}
le thread pendant les attentes I/O (requetes reseau, base de donnees, disque).
Pendant qu'une requete attend une reponse externe, l'event loop peut traiter
d'autres requetes en parallele logique. Pour une API qui passe le plus clair
de son temps a attendre la base de donnees ou un service externe, cela permet
de servir beaucoup plus de requetes concurrentes avec les memes ressources
qu'en mode synchrone bloquant.
\end{qa}

\begin{qa}{3}{Quel est le piege classique du code async : melanger appel bloquant et async ?}
Appeler une fonction \emph{synchrone bloquante} (ex. \texttt{requests.get()},
un calcul CPU lourd, un driver DB non-async) depuis une coroutine bloque tout
l'event loop -- toutes les autres taches concurrentes sont gelees pendant ce
temps, annulant le benefice de l'asynchrone. Il faut soit utiliser une
librairie async-native (\texttt{httpx}, \texttt{asyncpg}), soit deporter
l'appel bloquant dans un threadpool (\texttt{run\_in\_executor}, ou
\texttt{fastapi}'s \texttt{run\_in\_threadpool}).
\begin{piege}
Un \texttt{def} classique appele sans \texttt{await} dans un contexte async
n'est \textbf{pas} execute en parallele : il bloque comme du code synchrone
normal.
\end{piege}
\end{qa}

\newpage
% =============================================================================
\section{Django}
% =============================================================================

\subsection{Architecture}

\begin{qa}{1}{Explique le modele MTV de Django.}
Model--Template--View, variante du MVC. Le \textbf{Model} definit la structure
des donnees et la logique d'acces (ORM). Le \textbf{Template} gere le rendu
HTML (presentation). La \textbf{View} contient la logique metier : elle recoit
la requete, interroge les Models, choisit et alimente le Template. Django lui
meme joue le role du \og{}Controller\fg{} du MVC classique (routage, dispatch).
\end{qa}

\begin{qa}{2}{Que se passe-t-il lorsqu'une requete HTTP arrive dans Django ?}
1) Le serveur WSGI/ASGI recoit la requete et la transmet a Django.
2) Les \textbf{middlewares} la traitent dans l'ordre (auth, session, CSRF...).
3) Le \textbf{resolveur d'URL} (\texttt{urls.py}) trouve la vue correspondante.
4) La \textbf{vue} s'execute : elle interroge les \textbf{models} via l'ORM,
applique la logique metier.
5) La vue renvoie une \texttt{HttpResponse} (souvent generee via un
\textbf{template}, ou du JSON avec DRF).
6) La reponse retraverse les middlewares (dans l'ordre inverse) avant d'etre
renvoyee au client.
\end{qa}

\subsection{Routing, vues, middleware}

\begin{qa}{1}{A quoi servent \texttt{urls.py}, \texttt{path()} et \texttt{include()} ?}
\texttt{urls.py} declare la table de routage de l'application : association
entre motifs d'URL et vues. \texttt{path()} definit une route unique (avec
parametres typables, ex. \texttt{<int:id>}). \texttt{include()} permet de
deleguer un prefixe d'URL vers le \texttt{urls.py} d'une autre application,
favorisant la modularite (chaque app Django gere ses propres routes).
\end{qa}

\begin{qa}{1}{Function Based View vs Class Based View : pourquoi utiliser une CBV ?}
Une FBV est une simple fonction qui recoit la requete ; une CBV encapsule le
comportement dans une classe (souvent en heritant de vues generiques comme
\texttt{ListView}, \texttt{CreateView}). Les CBV favorisent la reutilisation
via l'heritage et les mixins, reduisent le code repetitif pour les patterns
CRUD standards, et separent proprement les methodes par verbe HTTP
(\texttt{get()}, \texttt{post()}...). Les FBV restent plus lisibles pour une
logique simple ou tres specifique.
\end{qa}

\begin{qa}{2}{Qu'est-ce qu'un middleware Django ?}
Un composant qui s'insere dans le traitement de \emph{chaque} requete/reponse,
sous forme de couches successives (pipeline). Utilise pour des
preoccupations transverses : authentification de session, gestion CSRF,
compression, logging, gestion des CORS, limitation de debit. Chaque
middleware peut modifier la requete avant qu'elle atteigne la vue, ou la
reponse avant qu'elle reparte vers le client.
\begin{lstlisting}
class TempsReponseMiddleware:
    def __init__(self, get_response):
        self.get_response = get_response

    def __call__(self, request):
        debut = time.time()
        response = self.get_response(request)
        response["X-Temps-Reponse"] = f"{time.time() - debut:.3f}"
        return response
\end{lstlisting}
\end{qa}

\subsection{ORM}

\begin{qa}{1}{Difference entre \texttt{filter()} et \texttt{get()} ?}
\texttt{filter()} renvoie toujours un \texttt{QuerySet} (potentiellement
vide, avec 0, 1 ou plusieurs resultats), evalue de maniere paresseuse.
\texttt{get()} renvoie \emph{un seul objet} et leve
\texttt{DoesNotExist} si aucun resultat, ou
\texttt{MultipleObjectsReturned} si plusieurs correspondent -- a utiliser
uniquement quand on est certain qu'au plus un enregistrement correspond (ex.
recherche par cle primaire).
\end{qa}

\begin{qa}{1}{Que fait \texttt{exclude()} ?}
L'inverse de \texttt{filter()} : renvoie un \texttt{QuerySet} contenant les
objets qui \emph{ne} correspondent \emph{pas} a la condition donnee.
\end{qa}

\begin{qa}{2}{Que sont \texttt{annotate()} et \texttt{aggregate()} ?}
\texttt{aggregate()} calcule une valeur agregee sur l'ensemble du queryset
(ex. \texttt{Avg}, \texttt{Count}, \texttt{Sum}) et renvoie un dictionnaire
unique. \texttt{annotate()} calcule une valeur agregee \emph{par ligne}
(souvent apres un \texttt{group by} implicite), en ajoutant un champ calcule a
chaque objet du queryset -- par exemple, le nombre de commandes par client.
\begin{lstlisting}
from django.db.models import Count, Avg

Client.objects.annotate(nb_commandes=Count("commandes"))
Commande.objects.aggregate(total=Avg("montant"))
\end{lstlisting}
\end{qa}

\begin{qa}{2}{Que sont \texttt{select\_related()} et \texttt{prefetch\_related()} ? Pourquoi sont-ils critiques pour la performance ?}
Ce sont des optimisations pour eviter le probleme des \textbf{requetes N+1}.
\texttt{select\_related()} effectue une jointure SQL (\texttt{JOIN}) pour
recuperer en une seule requete les objets lies via \texttt{ForeignKey} ou
\texttt{OneToOneField} (relations \og{}vers un seul\fg{}).
\texttt{prefetch\_related()} effectue une requete SQL \emph{separee}
supplementaire puis fait la jointure en Python -- necessaire pour les relations
\texttt{ManyToManyField} ou les relations inverses (\og{}vers plusieurs\fg{}),
ou un simple \texttt{JOIN} SQL ne suffit pas.
\begin{lstlisting}
# N+1 : 1 requete + 1 requete par article pour son auteur
articles = Article.objects.all()
for a in articles:
    print(a.auteur.nom)          # requete SQL a chaque iteration !

# Corrige : 1 seule requete grace au JOIN
articles = Article.objects.select_related("auteur")
\end{lstlisting}
\end{qa}

\begin{qa}{2}{Qu'est-ce que le probleme N+1 et comment le detecter ?}
Il survient quand on execute 1 requete pour recuperer une liste d'objets, puis
1 requete \emph{supplementaire par objet} pour recuperer une donnee liee
(ex. l'auteur de chaque article dans une boucle) -- soit $N+1$ requetes au
lieu d'une seule optimisee. On le detecte via \texttt{django-debug-toolbar}
(compteur de requetes SQL), les logs SQL en developpement, ou des outils
d'APM en production (New Relic, Datadog). Se corrige avec
\texttt{select\_related}/\texttt{prefetch\_related}.
\end{qa}

\subsection{Relations, migrations, admin, signals, forms}

\begin{qa}{1}{Difference entre \texttt{ForeignKey}, \texttt{OneToOneField} et \texttt{ManyToManyField} ?}
\texttt{ForeignKey} : relation plusieurs-vers-un (plusieurs commandes pour un
client). \texttt{OneToOneField} : relation un-vers-un stricte (un utilisateur
et son profil etendu). \texttt{ManyToManyField} : relation plusieurs-vers-plusieurs,
implementee via une table intermediaire (ex. articles et tags).
\end{qa}

\begin{qa}{1}{Pourquoi \texttt{makemigrations} et \texttt{migrate} sont-ils deux commandes distinctes ?}
\texttt{makemigrations} \emph{genere} les fichiers de migration en comparant
l'etat des models a la derniere migration connue (c'est une etape locale,
relisable et versionnable dans git). \texttt{migrate} \emph{applique}
reellement ces migrations a la base de donnees. Separer les deux permet de
relire/modifier une migration avant de l'executer, et de deployer la meme
migration de maniere identique sur plusieurs environnements (dev, staging,
prod) sans la regenerer a chaque fois.
\end{qa}

\begin{qa}{1}{Pourquoi Django Admin existe-t-il ?}
Interface d'administration generee automatiquement a partir des models,
permettant de gerer les donnees (CRUD) sans ecrire de vues ni de templates.
Tres utile en developpement et pour donner un back-office rapide a des
utilisateurs non techniques, mais generalement pas destine a etre expose
publiquement tel quel en production sans durcissement (permissions, 2FA, IP
whitelisting).
\end{qa}

\begin{qa}{2}{Qu'est-ce qu'un signal Django (\texttt{post\_save}, \texttt{pre\_save}) et quand l'utiliser ?}
Un signal permet a du code decouple de reagir a un evenement du cycle de vie
d'un model (avant/apres sauvegarde, suppression...) sans modifier directement
la methode \texttt{save()}. Exemple : creer automatiquement un
\texttt{Profile} a chaque creation d'\texttt{User} via \texttt{post\_save}. A
utiliser avec parcimonie : les signaux rendent le flux d'execution moins
explicite (\og{}action a distance\fg{}) et peuvent compliquer le debogage --
souvent, surcharger \texttt{save()} ou utiliser un service explicite est plus
lisible.
\end{qa}

\begin{qa}{2}{A quoi sert \texttt{ModelForm} et la methode \texttt{clean()} ?}
\texttt{ModelForm} genere automatiquement un formulaire (champs, widgets,
validations de base) a partir d'un \texttt{Model}, evitant la duplication
entre schema de donnees et schema de formulaire. La methode \texttt{clean()}
(ou \texttt{clean\_<champ>()}) permet d'ajouter de la validation
personnalisee, executee apres les validations de champ individuelles,
notamment pour verifier la coherence entre plusieurs champs.
\end{qa}

\newpage
% =============================================================================
\section{Django REST Framework (DRF)}
% =============================================================================

\begin{qa}{1}{Pourquoi utiliser DRF plutot que Django seul pour une API ?}
Django seul ne fournit pas nativement les briques attendues d'une API REST
moderne : serialisation/deserialisation JSON propre, negociation de contenu,
authentification par token, pagination, permissions granulaires,
documentation OpenAPI. DRF fournit ces briques de maniere standardisee et
extensible, tout en restant integre a l'ORM et aux permissions Django
existantes.
\end{qa}

\begin{qa}{1}{Difference entre \texttt{Serializer} et \texttt{ModelSerializer} ?}
\texttt{Serializer} se definit champ par champ manuellement (controle total,
utile pour des structures qui ne correspondent a aucun model). Un
\texttt{ModelSerializer} genere automatiquement les champs et validateurs a
partir d'un \texttt{Model} Django (comme \texttt{ModelForm} pour les
formulaires), reduisant fortement le code repetitif pour les cas standards.
\end{qa}

\begin{qa}{2}{Difference entre \texttt{APIView}, \texttt{GenericAPIView} et \texttt{ViewSet} ?}
\texttt{APIView} est le niveau le plus bas : on ecrit soi-meme
\texttt{get()}/\texttt{post()}, controle total mais plus verbeux.
\texttt{GenericAPIView} ajoute des comportements standards (queryset,
serializer\_class, pagination) combinables avec des \emph{mixins}
(\texttt{ListModelMixin}, \texttt{CreateModelMixin}...) pour construire du
CRUD avec peu de code. \texttt{ViewSet} regroupe la logique CRUD complete
d'une ressource dans une seule classe (\texttt{list}, \texttt{retrieve},
\texttt{create}, \texttt{update}, \texttt{destroy}), pensee pour etre branchee
a un \texttt{Router} qui genere automatiquement les URLs correspondantes.
\end{qa}

\begin{qa}{2}{Role d'un \texttt{Router} en DRF ?}
Genere automatiquement les routes URL standards pour un \texttt{ViewSet} (ex.
\texttt{GET /users/}, \texttt{POST /users/}, \texttt{GET /users/\{id\}/}...)
sans avoir a les declarer manuellement dans \texttt{urls.py}, en suivant les
conventions REST.
\end{qa}

\begin{qa}{1}{Comment fonctionnent \texttt{Permission} et \texttt{Authentication} en DRF ?}
\textbf{Authentication} determine \emph{qui} fait la requete (identification :
session, token, JWT...). \textbf{Permission} determine \emph{ce que
l'utilisateur identifie a le droit de faire} (autorisation : lecture seule,
proprietaire uniquement, admin...). Les deux s'executent avant que la vue ne
traite la requete, et peuvent etre combinees/personnalisees par vue ou
globalement.
\end{qa}

\begin{qa}{2}{Pourquoi la pagination est-elle importante sur une API qui liste des ressources ?}
Sans pagination, un endpoint de liste renverrait potentiellement des millions
d'enregistrements en une seule reponse, saturant la memoire serveur, le
reseau et le client. La pagination (par page, par curseur, ou par
offset/limit) borne la taille de chaque reponse et permet au client de
parcourir les resultats progressivement. DRF fournit
\texttt{PageNumberPagination}, \texttt{LimitOffsetPagination} et
\texttt{CursorPagination} (plus stable sur des donnees qui changent).
\end{qa}

\begin{qa}{2}{Difference entre filtering et ordering dans une API DRF ?}
Le \emph{filtering} restreint les resultats selon des criteres
(\texttt{?statut=actif}), tandis que l'\emph{ordering} controle l'ordre de
tri des resultats (\texttt{?ordering=-date\_creation}). Les deux sont souvent
geres via des parametres de query string et des backends dedies
(\texttt{django-filter}, \texttt{OrderingFilter}) plutot que codes en dur dans
chaque vue.
\end{qa}

\newpage
% =============================================================================
\section{FastAPI}
% =============================================================================

\begin{qa}{1}{Qu'est-ce que FastAPI et pourquoi est-il reconnu pour sa rapidite ?}
Un framework web Python moderne, base sur \textbf{Starlette} (couche ASGI) et
\textbf{Pydantic} (validation de donnees), concu nativement pour l'asynchrone.
Sa rapidite d'execution vient de son fondement ASGI/async (haute concurrence
I/O) et de Pydantic v2 (validation ecrite en Rust). Sa rapidite de
\emph{developpement} vient de la generation automatique de documentation
OpenAPI et de la validation automatique des types via les annotations Python.
\end{qa}

\begin{qa}{1}{Quelle est la difference entre WSGI et ASGI ?}
WSGI (\emph{Web Server Gateway Interface}) est synchrone : chaque requete
occupe un worker jusqu'a sa fin, ce qui est inefficace pour les operations
longues (I/O). ASGI (\emph{Asynchronous Server Gateway Interface}) supporte
nativement l'asynchrone, le WebSocket et les connexions longue duree,
permettant a un seul worker de gerer de nombreuses requetes concurrentes
pendant qu'elles attendent des I/O. Django classique tourne en WSGI (Gunicorn),
FastAPI tourne en ASGI (Uvicorn).
\end{qa}

\begin{qa}{2}{Pourquoi definir une route avec \texttt{async def} plutot que \texttt{def} simple dans FastAPI ?}
Avec \texttt{async def}, la fonction s'execute directement dans l'event loop :
ideal si tout le corps de la fonction est non-bloquant (appels a des
librairies async comme \texttt{asyncpg}, \texttt{httpx}). Avec un simple
\texttt{def}, FastAPI execute automatiquement la fonction dans un
\emph{threadpool} separe pour ne pas bloquer l'event loop -- utile si le code
appelle des librairies synchrones bloquantes. Le piege est d'utiliser
\texttt{async def} tout en appelant du code bloquant a l'interieur : cela
gele l'event loop entier.
\end{qa}

\subsection{Pydantic}

\begin{qa}{1}{A quoi sert \texttt{BaseModel} de Pydantic ?}
Classe de base pour definir des schemas de donnees types : Pydantic valide
automatiquement les types (et leve une erreur claire sinon), convertit/parse
les donnees brutes (JSON, dict) vers des objets Python types, et permet la
serialisation inverse (objet $\rightarrow$ dict/JSON). C'est le mecanisme qui
alimente la validation automatique des requetes et la generation du schema
OpenAPI dans FastAPI.
\begin{lstlisting}
from pydantic import BaseModel

class UtilisateurCreate(BaseModel):
    nom: str
    email: str
    age: int | None = None
\end{lstlisting}
\end{qa}

\begin{qa}{2}{Difference entre validation et serialisation dans Pydantic ?}
La \textbf{validation} (parsing) convertit et verifie des donnees brutes
entrantes (dict, JSON) vers une instance typee du modele, en levant une erreur
detaillee si un champ est invalide ou manquant. La \textbf{serialisation} fait
l'operation inverse : transformer l'instance du modele en dict/JSON pour la
renvoyer au client, avec un controle fin (exclusion de champs, alias, mode
JSON vs Python).
\end{qa}

\subsection{Dependances, requetes, reponses}

\begin{qa}{1}{Que fait \texttt{Depends()} et pourquoi l'utiliser ?}
Mecanisme d'\emph{injection de dependances} de FastAPI : on declare une
fonction (ou classe) qui produit une valeur (connexion DB, utilisateur
authentifie, configuration) et FastAPI l'execute automatiquement avant la
route, en injectant le resultat comme parametre. Cela permet de
\textbf{reutiliser} de la logique commune (auth, pagination, session DB) entre
plusieurs endpoints sans duplication, tout en restant facilement testable
(on peut surcharger une dependance dans les tests).
\begin{lstlisting}
def get_db():
    db = SessionLocal()
    try:
        yield db
    finally:
        db.close()

@app.get("/users/{id}")
def get_user(id: int, db: Session = Depends(get_db)):
    return db.query(User).get(id)
\end{lstlisting}
\end{qa}

\begin{qa}{1}{Comment recuperer Path Parameters, Query Parameters, Headers, Cookies et Body dans FastAPI ?}
Il suffit d'annoter les parametres de la fonction de route : un parametre
present dans le chemin de l'URL est un \emph{path parameter} automatique ; un
parametre simple non present dans le chemin devient un \emph{query
parameter} ; \texttt{Header()}, \texttt{Cookie()} et \texttt{Body()} (ou un
\texttt{BaseModel}) rendent explicite la source des autres donnees. FastAPI
deduit la source depuis la signature de la fonction et le type declare.
\begin{lstlisting}
@app.get("/items/{item_id}")
def read_item(item_id: int, q: str | None = None,
              x_token: str = Header(...)):
    ...
\end{lstlisting}
\end{qa}

\begin{qa}{2}{Quelles sont les principales classes de reponse (\texttt{Response}) en FastAPI ?}
\texttt{JSONResponse} (reponse JSON explicite, utile pour un code de statut ou
des headers personnalises), \texttt{Response} (reponse brute generique),
\texttt{StreamingResponse} (envoi progressif de donnees, utile pour de gros
fichiers ou du streaming de reponse LLM), \texttt{FileResponse} (envoi
efficace d'un fichier depuis le disque, avec support des headers de cache et
range requests).
\end{qa}

\begin{qa}{2}{Comment ajouter un middleware dans FastAPI ?}
Via le decorateur \texttt{@app.middleware("http")} ou
\texttt{app.add\_middleware(...)}. Un middleware HTTP recoit la requete,
peut la modifier ou executer du code avant, appelle
\texttt{await call\_next(request)} pour continuer le traitement, puis peut
modifier la reponse avant de la renvoyer.
\begin{lstlisting}
@app.middleware("http")
async def ajouter_temps_traitement(request, call_next):
    debut = time.time()
    response = await call_next(request)
    response.headers["X-Temps"] = str(time.time() - debut)
    return response
\end{lstlisting}
\end{qa}

\begin{qa}{2}{Pourquoi utiliser \texttt{BackgroundTasks} ?}
Pour executer une operation \emph{apres} avoir renvoye la reponse au client,
sans le faire attendre (ex. envoi d'un email de confirmation, ecriture de
log, notification). Contrairement a une vraie file de taches (Celery, RQ),
\texttt{BackgroundTasks} s'execute dans le meme processus juste apres la
reponse : adapte a des taches courtes et non critiques, pas a des traitements
lourds ou qui necessitent des garanties de fiabilite (retry, persistance).
\end{qa}

\begin{qa}{2}{A quoi servent les \emph{lifespan events} (\texttt{startup}/\texttt{shutdown}) ?}
Permettent d'executer du code une seule fois au demarrage de l'application
(ex. ouvrir un pool de connexions DB, charger un modele ML en memoire) et au
moment de l'arret (fermer proprement les connexions, liberer les ressources).
Depuis FastAPI recent, on utilise le context manager \texttt{lifespan} plutot
que les evenements \texttt{@app.on\_event} (deprecies).
\begin{lstlisting}
@asynccontextmanager
async def lifespan(app: FastAPI):
    await connect_to_db()      # startup
    yield
    await disconnect_from_db()  # shutdown

app = FastAPI(lifespan=lifespan)
\end{lstlisting}
\end{qa}

\begin{qa}{1}{Pourquoi la documentation Swagger/OpenAPI est-elle generee automatiquement par FastAPI ?}
Parce que FastAPI construit le schema OpenAPI directement a partir des
\emph{annotations de type} des fonctions de route et des modeles Pydantic --
il n'y a pas de documentation a maintenir separement du code : le schema
(donc la doc \texttt{/docs} via Swagger UI et \texttt{/redoc}) reste
automatiquement synchronise avec l'implementation reelle.
\end{qa}

\newpage
% =============================================================================
\section{REST API}
% =============================================================================

\begin{qa}{1}{Rappelle le role des principaux verbes HTTP.}
\texttt{GET} : lire une ressource (sans effet de bord). \texttt{POST} : creer
une nouvelle ressource / declencher une action. \texttt{PUT} : remplacer
integralement une ressource existante. \texttt{PATCH} : modifier partiellement
une ressource. \texttt{DELETE} : supprimer une ressource.
\end{qa}

\begin{qa}{1}{Quelle difference entre \texttt{PUT} et \texttt{PATCH} ?}
\texttt{PUT} remplace la ressource \emph{entiere} : le client doit envoyer
toutes les donnees, les champs omis sont generalement reinitialises.
\texttt{PATCH} applique une modification \emph{partielle} : seuls les champs
envoyes sont mis a jour, les autres restent inchanges.
\end{qa}

\begin{qa}{2}{Qu'est-ce que l'idempotence en REST ? Quelles methodes sont idempotentes ?}
Une operation est idempotente si l'executer plusieurs fois produit le meme
resultat/etat final que l'executer une seule fois (le client peut donc la
reessayer sans risque apres un timeout). \texttt{GET}, \texttt{PUT},
\texttt{DELETE} sont idempotentes ; \texttt{POST} ne l'est generalement pas
(appeler deux fois \texttt{POST /commandes} cree deux commandes). Notion
importante pour concevoir des retries reseau surs.
\end{qa}

\begin{qa}{1}{Explique les codes de statut HTTP les plus courants.}
\textbf{2xx (succes)} : \texttt{200} OK, \texttt{201} Created (apres un
\texttt{POST} reussi), \texttt{204} No Content (succes sans corps de reponse,
souvent apres \texttt{DELETE}). \textbf{4xx (erreur client)} : \texttt{400}
Bad Request (requete malformee), \texttt{401} Unauthorized (non authentifie),
\texttt{403} Forbidden (authentifie mais non autorise), \texttt{404} Not
Found, \texttt{409} Conflict (etat incompatible, ex. doublon), \texttt{422}
Unprocessable Entity (validation echouee -- tres frequent avec FastAPI/Pydantic).
\textbf{5xx (erreur serveur)} : \texttt{500} Internal Server Error.
\end{qa}

\newpage
% =============================================================================
\section{SQL}
% =============================================================================

\begin{qa}{1}{Difference entre \texttt{INNER JOIN} et \texttt{LEFT JOIN} ?}
\texttt{INNER JOIN} ne renvoie que les lignes ayant une correspondance dans
\emph{les deux} tables. \texttt{LEFT JOIN} renvoie toutes les lignes de la
table de gauche, meme sans correspondance a droite (les colonnes de droite
sont alors \texttt{NULL}) -- utile pour lister par exemple tous les clients,
y compris ceux sans commande.
\end{qa}

\begin{qa}{1}{Difference entre \texttt{WHERE} et \texttt{HAVING} (avec \texttt{GROUP BY}) ?}
\texttt{WHERE} filtre les lignes \emph{avant} l'agregation (\texttt{GROUP BY}),
ligne par ligne, et ne peut pas referencer une fonction d'agregation.
\texttt{HAVING} filtre les groupes \emph{apres} l'agregation, et peut donc
utiliser des fonctions comme \texttt{COUNT()}, \texttt{SUM()}.
\begin{lstlisting}
SELECT client_id, COUNT(*) AS nb
FROM commandes
WHERE statut = 'validee'
GROUP BY client_id
HAVING COUNT(*) > 5;
\end{lstlisting}
\end{qa}

\begin{qa}{2}{A quoi sert un \texttt{INDEX} et quel est son cout ?}
Un index accelere les recherches/filtres/tris (\texttt{WHERE},
\texttt{ORDER BY}, \texttt{JOIN}) sur une colonne, en evitant un scan complet
de la table (structure typiquement en B-tree, en \texttt{O(log n)} plutot que
\texttt{O(n)}). Le cout : chaque \texttt{INSERT}/\texttt{UPDATE}/\texttt{DELETE}
doit aussi mettre a jour l'index, ce qui ralentit les ecritures et consomme de
l'espace disque supplementaire -- il faut donc indexer les colonnes
frequemment filtrees, pas toutes les colonnes.
\end{qa}

\begin{qa}{1}{Difference entre \texttt{PRIMARY KEY}, \texttt{FOREIGN KEY} et \texttt{UNIQUE} ?}
\texttt{PRIMARY KEY} identifie de maniere unique chaque ligne d'une table
(implique unicite + non-nullite, indexee automatiquement).
\texttt{FOREIGN KEY} garantit l'integrite referentielle en liant une colonne
a la cle primaire d'une autre table. \texttt{UNIQUE} garantit qu'une colonne
(ou combinaison de colonnes) ne contient pas de doublon, sans en faire
l'identifiant de la ligne.
\end{qa}

\begin{qa}{2}{Qu'est-ce qu'une transaction SQL ?}
Un ensemble d'operations executees comme une unite atomique : soit toutes
reussissent (\texttt{COMMIT}), soit aucune n'est appliquee
(\texttt{ROLLBACK} en cas d'erreur). Essentiel pour garantir la coherence des
donnees lors d'operations multi-etapes (ex. debiter un compte et en crediter
un autre doivent reussir ensemble ou pas du tout).
\end{qa}

\begin{qa}{2}{Exercice -- Ecris une requete qui recupere les 10 derniers utilisateurs inscrits.}
\begin{lstlisting}
SELECT *
FROM utilisateurs
ORDER BY date_inscription DESC
LIMIT 10;
\end{lstlisting}
\end{qa}

\begin{qa}{2}{Exercice -- Compte le nombre d'utilisateurs par pays, tries par nombre decroissant.}
\begin{lstlisting}
SELECT pays, COUNT(*) AS nb_utilisateurs
FROM utilisateurs
GROUP BY pays
ORDER BY nb_utilisateurs DESC;
\end{lstlisting}
\end{qa}

\newpage
% =============================================================================
\section{PostgreSQL}
% =============================================================================

\begin{qa}{2}{Qu'est-ce que ACID ?}
Quatre garanties attendues d'une transaction fiable : \textbf{Atomicite}
(tout ou rien), \textbf{Coherence} (la base passe d'un etat valide a un autre
etat valide, respecte les contraintes), \textbf{Isolation} (des transactions
concurrentes ne s'interferent pas de maniere incoherente), \textbf{Durabilite}
(une fois validee, une transaction survit meme a une panne).
\end{qa}

\begin{qa}{3}{Que sont les niveaux d'isolation (Isolation Level) ?}
Ils definissent a quel point une transaction est protegee des effets des
autres transactions concurrentes, du plus permissif au plus strict :
\texttt{READ UNCOMMITTED}, \texttt{READ COMMITTED} (defaut PostgreSQL),
\texttt{REPEATABLE READ}, \texttt{SERIALIZABLE}. Plus le niveau est strict,
plus les anomalies (dirty read, non-repeatable read, phantom read) sont
evitees, mais plus le risque de contention/rollback augmente.
\end{qa}

\begin{qa}{3}{Qu'est-ce qu'un deadlock et comment l'eviter ?}
Situation ou deux (ou plus) transactions s'attendent mutuellement pour
liberer des verrous qu'elles detiennent, bloquant indefiniment -- PostgreSQL
detecte ces cycles et annule automatiquement une des transactions
(erreur \texttt{deadlock detected}). Pour limiter le risque : toujours
acquerir les verrous/mettre a jour les lignes dans le \emph{meme ordre} dans
toute l'application, garder les transactions courtes, eviter les
transactions interactives qui attendent une action utilisateur en cours de
route.
\end{qa}

\newpage
% =============================================================================
\section{Securite}
% =============================================================================

\begin{qa}{1}{Difference entre Authentication et Authorization ?}
\textbf{Authentication} (authentification) repond a \og{}qui es-tu ?\fg{}
(verifier l'identite : mot de passe, token). \textbf{Authorization}
(autorisation) repond a \og{}as-tu le droit de faire ceci ?\fg{} (verifier les
permissions une fois l'identite etablie).
\end{qa}

\begin{qa}{2}{Qu'est-ce qu'un JWT et comment fonctionne-t-il ?}
JSON Web Token : un jeton auto-suffisant compose de trois parties encodees en
base64 et separees par des points -- \emph{header} (algorithme),
\emph{payload} (donnees/claims, ex. id utilisateur, expiration),
\emph{signature} (garantit l'integrite, calculee avec une cle secrete ou une
paire de cles). Le serveur peut verifier la validite d'un JWT sans requete
base de donnees (juste verifier la signature), ce qui le rend adapte aux
architectures distribuees/stateless.
\begin{piege}
Un JWT n'est pas chiffre par defaut (juste signe) : ne jamais y mettre de
donnees sensibles en clair, tout le monde peut decoder le payload.
\end{piege}
\end{qa}

\begin{qa}{2}{Difference entre Access Token et Refresh Token ?}
L'\textbf{access token} a une duree de vie courte et est envoye a chaque
requete pour prouver l'identite -- s'il est vole, l'exposition est limitee
dans le temps. Le \textbf{refresh token} a une duree de vie plus longue et
sert uniquement a obtenir un nouvel access token sans redemander les
identifiants, generalement stocke de maniere plus securisee (cookie
httpOnly) et revocable cote serveur.
\end{qa}

\begin{qa}{2}{Explique le principe general d'OAuth2.}
Protocole d'autorisation delegue : un utilisateur autorise une application
tierce a acceder a ses ressources (sur un autre service) sans lui partager
son mot de passe. L'application recoit un \emph{access token} avec un scope
limite, delivre par un serveur d'autorisation apres consentement explicite de
l'utilisateur. FastAPI et DRF fournissent tous deux des integrations pour les
flows OAuth2 standards (ex. \emph{password flow}, \emph{authorization code}).
\end{qa}

\begin{qa}{1}{Pourquoi ne jamais stocker un mot de passe en clair ? Role de \texttt{bcrypt}.}
Si la base de donnees fuite, des mots de passe en clair exposent
immediatement tous les comptes (et souvent d'autres services, via la
reutilisation de mots de passe). On stocke a la place un \emph{hash}
irreversible du mot de passe. \texttt{bcrypt} (comme \texttt{argon2}) est
concu pour etre volontairement \emph{lent} et inclut un \emph{salt} (valeur
aleatoire) unique par mot de passe, ce qui rend les attaques par force brute
et les rainbow tables couteuses, contrairement a un hash rapide generique
comme MD5 ou SHA-256 seul.
\end{qa}

\begin{qa}{2}{Qu'est-ce que CORS et pourquoi le navigateur le bloque-t-il par defaut ?}
\emph{Cross-Origin Resource Sharing} : mecanisme de securite du navigateur qui
bloque par defaut les requetes JavaScript vers un domaine different de celui
qui a servi la page (\emph{same-origin policy}), pour empecher un site
malveillant de faire des requetes authentifiees vers un autre service au nom
de l'utilisateur. Le serveur doit explicitement autoriser les origines
souhaitees via des en-tetes \texttt{Access-Control-Allow-Origin}.
\end{qa}

\begin{qa}{2}{Qu'est-ce que CSRF et comment Django s'en protege-t-il ?}
\emph{Cross-Site Request Forgery} : une attaque ou un site malveillant fait
executer, a l'insu de l'utilisateur, une requete d'etat (ex. changement de
mot de passe) vers un site ou il est deja authentifie (via ses cookies de
session). Django genere un \emph{token CSRF} unique par session, inclus dans
chaque formulaire, et verifie sa presence/validite a chaque requete modifiant
l'etat (\texttt{POST}/\texttt{PUT}/\texttt{DELETE}) -- un attaquant externe ne
peut pas connaitre ce token.
\end{qa}

\begin{qa}{2}{Qu'est-ce que XSS et comment s'en proteger ?}
\emph{Cross-Site Scripting} : injection de code JavaScript malveillant dans
une page consultee par d'autres utilisateurs (via un champ de saisie non
echappe, par exemple). Protection : toujours \emph{echapper}/encoder les
donnees utilisateur avant de les injecter dans le HTML (les moteurs de
template Django echappent par defaut), valider les entrees, et definir une
politique \texttt{Content-Security-Policy} stricte.
\end{qa}

\begin{qa}{1}{Qu'est-ce qu'une injection SQL et comment l'ORM protege-t-il ?}
Une injection SQL survient quand une entree utilisateur est concatenee
directement dans une requete SQL, permettant a un attaquant d'y injecter du
SQL arbitraire (ex. \texttt{' OR '1'='1}). Les ORM (Django ORM, SQLAlchemy)
utilisent des \emph{requetes parametrees} par defaut (les valeurs sont
envoyees separement de la requete SQL, jamais concatenees textuellement), ce
qui neutralise ce risque tant qu'on evite le SQL brut non parametre
(\texttt{raw()}, \texttt{extra()}) avec des entrees utilisateur directes.
\end{qa}

\newpage
% =============================================================================
\section{Tests}
% =============================================================================

\begin{qa}{1}{Difference entre \texttt{pytest} et \texttt{unittest} ?}
\texttt{unittest} est le framework de test integre a la bibliotheque standard,
base sur des classes heritant de \texttt{TestCase} avec des methodes
\texttt{assertX}. \texttt{pytest} est une librairie tierce plus expressive :
simple \texttt{assert} natif, systeme de \emph{fixtures} puissant et
composable, tests parametrables, et un ecosysteme de plugins riche
(\texttt{pytest-django}, \texttt{pytest-asyncio}, \texttt{pytest-cov}). Les
deux restent compatibles entre eux.
\end{qa}

\begin{qa}{2}{A quoi sert une \emph{fixture} en pytest ?}
Une fonction reutilisable qui prepare un etat/une ressource necessaire a un
ou plusieurs tests (ex. une connexion DB de test, un utilisateur factice,
un client API), injectee automatiquement dans la fonction de test qui la
declare en parametre. Permet de factoriser la mise en place/le nettoyage
(\emph{setup}/\emph{teardown}) sans duplication.
\begin{lstlisting}
import pytest

@pytest.fixture
def utilisateur_test(db):
    return User.objects.create(username="alice")

def test_profil(utilisateur_test):
    assert utilisateur_test.username == "alice"
\end{lstlisting}
\end{qa}

\begin{qa}{2}{Quand et pourquoi utiliser \texttt{mock} dans un test ?}
Pour remplacer une dependance externe (appel reseau, service tiers, horloge
systeme, envoi d'email) par un objet controle, afin de rendre le test rapide,
deterministe et independant de systemes externes potentiellement instables
ou couteux. Le risque : sur-mocker peut faire passer un test qui ne reflete
plus le comportement reel du systeme integre.
\begin{lstlisting}
from unittest.mock import patch

@patch("app.services.envoyer_email")
def test_inscription_envoie_email(mock_email):
    inscrire_utilisateur("alice@test.com")
    mock_email.assert_called_once()
\end{lstlisting}
\end{qa}

\begin{qa}{1}{A quoi servent \texttt{APIClient} (DRF) et \texttt{TestClient} (FastAPI) ?}
Ce sont des clients HTTP factices permettant de tester les endpoints d'une
API sans lancer un vrai serveur reseau : ils appellent directement
l'application (WSGI/ASGI) en memoire, ce qui rend les tests rapides et
isoles. Ils permettent de simuler des requetes authentifiees, de verifier le
code de statut, le corps JSON et les en-tetes de la reponse.
\begin{lstlisting}
# FastAPI
from fastapi.testclient import TestClient
client = TestClient(app)
def test_health():
    assert client.get("/health").status_code == 200

# DRF
from rest_framework.test import APIClient
def test_liste_users():
    client = APIClient()
    response = client.get("/api/users/")
    assert response.status_code == 200
\end{lstlisting}
\end{qa}

\newpage
% =============================================================================
\section{Docker}
% =============================================================================

\begin{qa}{1}{Difference entre une \emph{image} et un \emph{container} Docker ?}
Une \textbf{image} est un modele en lecture seule (couches empilees) decrivant
un systeme de fichiers et une configuration d'execution. Un \textbf{container}
est une \emph{instance en cours d'execution} de cette image, avec sa propre
couche modifiable ephemere et son propre espace processus/reseau isole. Une
meme image peut donner naissance a plusieurs containers independants.
\end{qa}

\begin{qa}{1}{A quoi sert un \texttt{Dockerfile} ?}
Fichier texte decrivant, instruction par instruction, comment construire une
image Docker : image de base (\texttt{FROM}), copie de fichiers
(\texttt{COPY}), installation de dependances (\texttt{RUN}), variables
d'environnement (\texttt{ENV}), commande de demarrage (\texttt{CMD}).
\begin{lstlisting}
FROM python:3.12-slim
WORKDIR /app
COPY requirements.txt .
RUN pip install --no-cache-dir -r requirements.txt
COPY . .
CMD ["uvicorn", "main:app", "--host", "0.0.0.0", "--port", "8000"]
\end{lstlisting}
\end{qa}

\begin{qa}{2}{A quoi sert un \emph{volume} Docker ?}
Un mecanisme de persistance de donnees en dehors du cycle de vie ephemere
d'un container : les donnees ecrites dans un volume survivent a la
suppression du container (contrairement a la couche modifiable interne). Utile
pour les donnees de base de donnees, fichiers uploades, ou pour partager du
code entre l'hote et le container en developpement.
\end{qa}

\begin{qa}{2}{Pourquoi utiliser Docker Compose ?}
Pour definir et orchestrer plusieurs containers lies (ex. app +
PostgreSQL + Redis) via un seul fichier declaratif
(\texttt{docker-compose.yml}), avec leur reseau partage, leurs volumes et
leur ordre de demarrage, plutot que de lancer chaque container manuellement
avec de longues commandes \texttt{docker run}.
\end{qa}

\newpage
% =============================================================================
\section{Git}
% =============================================================================

\begin{qa}{1}{Difference entre \texttt{merge} et \texttt{rebase} ?}
\texttt{merge} combine deux branches en creant un \emph{commit de fusion}
qui preserve l'historique exact tel qu'il s'est produit (deux lignes qui se
rejoignent). \texttt{rebase} \emph{rejoue} les commits d'une branche par
dessus une autre, produisant un historique lineaire sans commit de fusion,
mais en \emph{reecrivant} les hash de commits -- a eviter sur une branche deja
partagee/pushee, car cela desynchronise les autres collaborateurs.
\end{qa}

\begin{qa}{2}{A quoi sert \texttt{git stash} ?}
Met de cote temporairement les modifications non commitees (working
directory + index) pour revenir a un etat propre, sans avoir a commiter du
travail inacheve -- utile pour changer rapidement de branche. On les
recupere ensuite avec \texttt{git stash pop} ou \texttt{git stash apply}.
\end{qa}

\begin{qa}{2}{A quoi sert \texttt{cherry-pick} ?}
Applique un commit specifique (identifie par son hash) d'une branche sur une
autre, sans fusionner tout l'historique des deux branches -- utile pour
recuperer un correctif precis sans embarquer d'autres changements non
souhaites.
\end{qa}

\begin{qa}{2}{Difference entre \texttt{git reset} et \texttt{git revert} ?}
\texttt{reset} deplace le pointeur de branche vers un commit anterieur,
\emph{reecrivant} l'historique (dangereux sur une branche partagee car les
commits \og{}annules\fg{} disparaissent de la branche). \texttt{revert} cree
un \emph{nouveau} commit qui annule les changements d'un commit precedent,
sans reecrire l'historique existant -- surer pour annuler un changement deja
pousse/partage.
\end{qa}

\newpage
% =============================================================================
\section{Architecture logicielle}
% =============================================================================

\begin{qa}{2}{Rappelle le principe SOLID.}
\textbf{S}ingle Responsibility -- une classe a une seule raison de changer.
\textbf{O}pen/Closed -- ouvert a l'extension, ferme a la modification.
\textbf{L}iskov Substitution -- une sous-classe doit pouvoir remplacer sa
classe mere sans casser le comportement attendu. \textbf{I}nterface
Segregation -- preferer plusieurs interfaces specifiques a une seule
interface generale. \textbf{D}ependency Inversion -- dependre d'abstractions,
pas d'implementations concretes.
\end{qa}

\begin{qa}{2}{Qu'est-ce que le Repository Pattern et quel probleme resout-il ?}
Un pattern qui isole la logique d'acces aux donnees (requetes DB) derriere
une interface dediee, decouplant la logique metier du mecanisme de
persistance concret (ORM, requetes SQL brutes, API externe). Avantages :
logique metier plus facile a tester (on peut injecter un faux repository en
test) et possibilite de changer de source de donnees sans toucher a la
logique metier.
\end{qa}

\begin{qa}{2}{Qu'est-ce qu'une \emph{Service Layer} et pourquoi la separer des vues ?}
Une couche intermediaire qui contient la logique metier (regles, orchestration
de plusieurs operations) independamment des vues/controleurs HTTP. Sans elle,
la logique metier finit eparpillee dans les vues, difficile a reutiliser
(ex. entre une route API et une tache planifiee) et a tester unitairement
sans monter tout le framework web.
\end{qa}

\begin{qa}{2}{Qu'est-ce que l'injection de dependances (Dependency Injection) et pourquoi FastAPI en fait un usage central ?}
Principe consistant a \emph{fournir} les dependances d'un composant depuis
l'exterieur (au lieu que le composant les construise lui-meme), ce qui
decouple le code et facilite les tests (substitution par un mock/faux
objet). \texttt{Depends()} de FastAPI est un systeme d'injection de
dependances natif : sessions DB, utilisateur authentifie, configuration --
tout peut etre injecte de maniere declarative et surchargee facilement en
test via \texttt{app.dependency\_overrides}.
\end{qa}

\newpage
% =============================================================================
\section{Performance}
% =============================================================================

\begin{qa}{2}{Comment reduire l'impact du probleme N+1 au-dela de \texttt{select\_related}/\texttt{prefetch\_related} ?}
En surveillant le nombre de requetes SQL par endpoint en continu (APM,
\texttt{django-debug-toolbar} en dev), en ecrivant des tests qui verifient un
nombre de requetes borne (\texttt{assertNumQueries}), et en revisant les
serializers DRF/Pydantic qui accedent souvent a des relations imbriquees sans
prefetch explicite.
\end{qa}

\begin{qa}{2}{Comment fonctionne le cache avec Redis dans une API Python ?}
Redis est une base cle-valeur en memoire, tres rapide, utilisee pour stocker
temporairement le resultat d'un calcul ou d'une requete couteuse (cache de
requete, cache de session, rate limiting, file de taches avec Celery). Le
pattern classique : verifier si la cle existe dans Redis (\emph{cache hit}),
sinon calculer/interroger la DB puis stocker le resultat dans Redis avec une
\emph{expiration} (\texttt{TTL}) avant de le renvoyer.
\end{qa}

\begin{qa}{2}{Difference entre pagination et \emph{lazy loading} ?}
La \textbf{pagination} decoupe explicitement les resultats en pages
demandees par le client (parametre \texttt{page}/\texttt{limit}). Le
\textbf{lazy loading} (chargement paresseux) retarde le chargement d'une
donnee jusqu'a ce qu'elle soit reellement necessaire (ex. un QuerySet Django
n'execute la requete SQL qu'au moment ou on itere dessus, pas a sa
declaration) -- les deux limitent le travail effectue en avance, mais a des
niveaux differents.
\end{qa}

\newpage
% =============================================================================
\section{Deploiement}
% =============================================================================

\begin{qa}{1}{Difference entre Gunicorn et Uvicorn ? Pourquoi parfois les deux ensemble ?}
\textbf{Gunicorn} est un serveur d'application WSGI, gerant un pool de
processus \emph{workers} (robustesse, gestion des crashs, redemarrage).
\textbf{Uvicorn} est un serveur ASGI performant capable de gerer l'asynchrone.
En production, on utilise souvent Gunicorn comme \emph{process manager}
(gestion des workers, redemarrage) avec des workers de type
\texttt{UvicornWorker}, combinant la robustesse operationnelle de Gunicorn et
les capacites async d'Uvicorn.
\end{qa}

\begin{qa}{2}{Pourquoi placer Nginx devant une application Django/FastAPI ?}
Nginx sert de \emph{reverse proxy} : il termine les connexions TLS/HTTPS,
sert les fichiers statiques efficacement (sans solliciter l'application
Python), gere la compression, le load balancing entre plusieurs workers/
instances, et protege l'application des connexions lentes/malveillantes
directes.
\end{qa}

\begin{qa}{2}{Pourquoi utiliser des variables d'environnement pour la configuration ?}
Pour separer la configuration (secrets, URLs de base de donnees, cles API)
du code source, permettant de deployer le \emph{meme} artefact/image sur
plusieurs environnements (dev, staging, prod) sans jamais commiter de secrets
dans le depot git -- principe des \emph{Twelve-Factor Apps}.
\end{qa}

\begin{qa}{2}{Que verifie typiquement un pipeline CI/CD pour une API Python ?}
Lint/formatage (ruff, black), typage statique (mypy), execution de la suite
de tests (pytest), build de l'image Docker, scan de securite des dependances,
puis deploiement automatique (souvent apres validation manuelle en
production). L'objectif est de detecter les regressions avant qu'elles
n'atteignent la production.
\end{qa}

\newpage
% =============================================================================
\section{Questions sur tes projets}
% =============================================================================
\begin{infobox}{Pourquoi cette section compte le plus}
30 a 40\,\% d'un entretien backend porte sur l'explication detaillee de tes
propres projets (CV, GitHub). Prepare des reponses concretes et chiffrees --
\og{}j'ai reduit le temps de reponse de 800ms a 120ms en ajoutant du cache
Redis\fg{} vaut toujours mieux qu'une reponse generique.
\end{infobox}

\begin{qa}{2}{Pourquoi Django plutot que FastAPI (ou l'inverse) pour ce projet ?}
Prepare une reponse basee sur les besoins reels du projet : Django s'impose
quand on a besoin d'un ecosysteme complet rapidement (admin, ORM mature,
auth, formulaires) pour une application classique avec beaucoup de CRUD et
de vues. FastAPI s'impose pour des APIs pures orientees performance/async,
des microservices, ou une forte contrainte de validation de schema stricte
(Pydantic) -- typiquement, une API consommee par un frontend decouple ou
d'autres services.
\end{qa}

\begin{qa}{2}{Pourquoi FastAPI plutot que Flask ?}
FastAPI apporte nativement : validation de types via Pydantic, documentation
OpenAPI automatique, support async natif (ASGI), et injection de dependances
integree. Flask est plus minimaliste (WSGI, synchrone par defaut), demandant
des extensions tierces pour obtenir des fonctionnalites equivalentes.
\end{qa}

\begin{qa}{2}{Comment as-tu structure ton projet (architecture des dossiers/couches) ?}
Prepare une reponse concrete : separation routes/schemas (Pydantic)/models
(ORM)/services (logique metier)/repositories (acces donnees), configuration
centralisee, tests miroir de la structure source. Sois pret a justifier
\emph{pourquoi} cette structure (testabilite, lisibilite, isolation des
couches) plutot que de la lister sans argumentation.
\end{qa}

\begin{qa}{2}{Comment geres-tu les exceptions dans ton API ?}
Prepare une reponse concrete : exceptions metier personnalisees, handler
global (\texttt{@app.exception\_handler} en FastAPI, middleware en Django)
qui traduit chaque exception en reponse HTTP coherente (code de statut +
message structure), sans jamais exposer de stack trace ou de details internes
au client en production.
\end{qa}

\begin{qa}{2}{Comment valides-tu les donnees entrantes ?}
Prepare une reponse concrete : Pydantic (FastAPI) ou Serializers/Forms
(Django/DRF) pour la validation de structure/type au plus tot (a la
frontiere de l'API), plus des regles metier explicites dans la couche
service pour les validations qui dependent de l'etat de la base (ex.
unicite, coherence entre entites).
\end{qa}

\begin{qa}{2}{Comment securises-tu ton API ?}
Prepare une reponse concrete couvrant : authentification (JWT/OAuth2),
autorisation par role/scope, validation stricte des entrees, HTTPS partout,
CORS configure explicitement (pas de wildcard en prod), rate limiting,
gestion des secrets via variables d'environnement (jamais en dur dans le
code).
\end{qa}

\begin{qa}{2}{Comment testes-tu ton code ?}
Prepare une reponse concrete : pyramide de tests (unitaires sur la logique
metier, tests d'integration sur les endpoints via TestClient/APIClient,
quelques tests end-to-end critiques), utilisation de fixtures/mocks pour
isoler les dependances externes, mesure de couverture pour reperer les zones
non testees sans en faire un objectif chiffre absolu.
\end{qa}

\begin{qa}{2}{Comment deploies-tu ton application ?}
Prepare une reponse concrete : conteneurisation Docker, pipeline CI/CD
(tests puis build puis deploiement), configuration via variables
d'environnement, reverse proxy (Nginx) devant Gunicorn/Uvicorn, strategie de
rollback en cas de probleme.
\end{qa}

\begin{qa}{2}{Pourquoi PostgreSQL plutot qu'un autre SGBD ?}
Prepare une reponse concrete : conformite ACID stricte, richesse des types
(JSONB pour des donnees semi-structurees, types tableaux), extensions
puissantes (pg\_vector pour l'IA/RAG, pg\_cron), excellent support des
transactions concurrentes et de l'indexation avancee, tres bon ecosysteme
avec les ORMs Python.
\end{qa}

\begin{qa}{2}{Pourquoi Redis dans ton architecture ?}
Prepare une reponse concrete : cas d'usage precis -- cache de requetes
couteuses, stockage de session, rate limiting, file de messages/taches
(avec Celery ou en broker pour Kafka-like patterns), pub/sub pour de la
notification temps reel. Justifie le choix par la latence (memoire, pas
disque) plutot que \og{}parce que c'est standard\fg{}.
\end{qa}

\begin{qa}{3}{Comment evites-tu les requetes N+1 dans ton projet concret ?}
Prepare une reponse concrete avec un exemple reel : identification via
profiling/logs SQL, correction avec \texttt{select\_related}/
\texttt{prefetch\_related} (Django) ou chargement explicite en batch
(SQLAlchemy \texttt{joinedload}/\texttt{selectinload}), et si possible un
chiffre avant/apres (nombre de requetes, temps de reponse).
\end{qa}

\begin{qa}{3}{Comment ameliores-tu les performances globales de ton API ?}
Prepare une reponse concrete et hierarchisee : d'abord mesurer (profiling,
APM) avant d'optimiser, puis leviers courants -- indexation SQL adaptee,
cache (Redis) sur les lectures couteuses et peu volatiles, pagination
systematique, requetes async pour les I/O concurrentes, mise en cache HTTP
(ETag, Cache-Control), et en dernier recours scaling horizontal.
\end{qa}

\end{document}
